\section{Introduction}
\label{sec:intro}

Developers regularly confront security decisions in the course of building and maintaining their own software. They need to know whether a freshly disclosed vulnerability in a dependency actually affects the way they use it, which release first introduced a regression, whether a proposed patch is sufficient, or which version is safe to upgrade to. When the answer is not obvious, they ask---on issue trackers, in advisory discussions, and on CVE-linked threads of the projects they depend on. How developers resolve such questions has direct consequences for the security of the software they ship: a substantial body of usable-security research shows that developers lean heavily on online information sources when solving security problems, and that the quality of those sources measurably shapes the security of the resulting code. Acar et al.\ found that only a small fraction of the Stack~Overflow threads developers consult for security tasks contain secure advice~\cite{acar2016information}, and Fischer et al.\ traced insecure snippets from such threads into hundreds of thousands of shipped applications~\cite{fischer2017stackoverflow}. Getting trustworthy answers to developer security questions is therefore not a convenience problem but a software-security problem.

These questions are also distinctive in what it takes to answer them. Unlike general code-comprehension questions, a developer security question typically cannot be resolved from the repository's source alone. Answering ``does CVE-2022-25883 affect me, and what do I upgrade to?'' requires cross-referencing artifacts that live partly outside the codebase: the CVE or GHSA advisory and its affected-version ranges, the fix commit or pull request, the project's dependency manifest, and sometimes an upstream RFC or registry. Project maintainers answer these questions precisely because they know which artifact to consult---and, crucially, \emph{different questions demand different artifacts}. A question about a transitive dependency is answered from the advisory and the manifest; a question about a TLS validation change is answered from the source and the commit history. Which artifact types are necessary for which kinds of security question is, to date, an unmeasured empirical question.

Large language models (LLMs) and coding agents are increasingly the first responders to such questions, yet no existing benchmark evaluates them on the security questions developers actually ask, or on the artifact dependence those questions exhibit. Repository-level QA benchmarks---SWE-QA~\cite{sweqa2026}, CoReQA~\cite{coreqa2025}, and StackRepoQA~\cite{stackrepoqa2025}---treat developer questions uniformly and consider only code-internal context; they neither isolate security questions nor model the external advisory artifacts those questions depend on. (StackRepoQA further reports that much of the apparent accuracy on such tasks reflects memorization of public threads rather than reasoning~\cite{stackrepoqa2025}, a contamination risk any security-QA benchmark must confront.) Security-specific LLM benchmarks, in turn, measure a different task: SecVulEval~\cite{secvuleval2025} and CodeSecEval~\cite{codeseceval2024} evaluate vulnerability detection and secure code generation, SecureAgentBench~\cite{secureagentbench2025} evaluates whether agents write secure patches, and SecQA~\cite{secqa2023} tests textbook security knowledge via multiple-choice questions. None of them studies real, free-text developer security questions, and none treats artifact context as a controlled variable.

We close this gap with \textbf{SecDevQA}, a benchmark of real developer security question--answer pairs mined from GitHub issues, discussions, and CVE-linked threads across open-source GitHub repositories. Two design choices make the benchmark both rigorous and analytically useful. First, we restrict the benchmark to \emph{hard-verifiable} pairs: pairs whose maintainer answer states at least one externally checkable fact---a CVE or GHSA identifier, a fixed version number, a fix commit or pull request, or an advisory URL. These hard facts anchor automated grading, turning the otherwise subjective task of scoring a free-text answer into the narrowly scoped task of checking whether a model's response correctly conveys a named, verifiable claim---a judgment a human can validate in seconds. Second, every pair is annotated with the artifact types its answer drew on, so that artifact context can be supplied to models as a controlled experimental condition rather than left implicit.

This design lets us pursue three questions that prior work cannot. We first derive, by inductive open coding of the mined corpus, a taxonomy of developer security question categories and an empirical map of which artifact types maintainers draw on to answer each---findings about developer security practice that stand independently of any model evaluation (RQ1). We then measure how accurately frontier LLMs, open-weight LLMs, and coding agents answer these questions, broken down by category (RQ2). Finally, and centrally, we ask which artifact types are \emph{necessary} to answer which categories of question, by supplying artifacts as context under controlled conditions and measuring each type's marginal contribution (RQ3); we further examine whether autonomous agents
retrieve the artifacts they actually need (RQ4). Our guiding hypothesis is that context dependence is category-specific: an artifact type that is essential for one category of security question yields little benefit for another---a result with direct implications for how retrieval-augmented and agentic security tools should be built.

\noindent\textbf{Contributions.} This paper makes the following contributions:
\begin{itemize}
  \item \textbf{SecDevQA}, a benchmark of real, hard-verifiable developer security Q\&A pairs mined from open-source repositories, each annotated with the artifact types needed to answer it and a structured record of its verifiable hard facts (\S\ref{sec:benchmark-construction}).
  \item A fact-grounded grading protocol that exploits hard facts to make LLM-judge scoring narrowly scoped and human-validatable
        (\S\ref{sec:benchmark-construction}).
  \item An inductively derived taxonomy of developer security question categories and an empirical category-to-artifact map, both independent of any LLM evaluation (\S\ref{sec:open-coding}).
  \item An evaluation of frontier and open-weight LLMs and coding agents under controlled context conditions, quantifying the marginal, category-specific contribution of each artifact type (\S\ref{sec:experimental-setup}).
  \item A public release of the benchmark, the mining and extraction pipeline, and the evaluation harness.
\end{itemize}
